% [原文 p.133]
\par\noindent（2）$\dim X=3$，$(X,0)$ 为终端 [Mor88]。

\par\noindent（3）$\dim X=3$，$(X,\Delta)$ 为 klt [Sho92]。新的证明以及向 $(X,\Delta)$ 为 lc 情形的一般化见 [K+92]。

第 7 章用于所谓三维半稳定翻转的存在证明。由三维翻转的分类 [KM92]，几乎所有的三维翻转都是半稳定的，但其证明并不容易。半稳定翻转的重要性源于其在曲面族上的应用。这在 §7.5 中说明。\hfill$\square$

二维情形不存在翻转，故此时 MMP 已完成。

\begin{theorem}{3.47}
设 $(X,\Delta)$ 为二维 Q-分解射影 dlt 对。则存在由至多 $\rho(X)-1$ 个收缩组成的列
\[
(X,\Delta)=(X_0,\Delta_0)\to(X_1,\Delta_1)\to\cdots\to(X_i,\Delta_i)=(X^{*},\Delta^{*}),
\]
满足下列条件之一。
\begin{enumerate}[label=(\arabic*)]
\item （对数极小模型（log minimal model））$K_{X^{*}}+\Delta^{*}$ 为 nef。
\item （对数直纹曲面（log ruled surface））存在到曲线 $C$ 的射 $g:X^{*}\to C$，使得 $-(K_{X^{*}}+\Delta^{*})$ 为 $g$-丰富。
\item （对数 Del Pezzo 曲面）$-(K_{X^{*}}+\Delta^{*})$ 为丰富。
\end{enumerate}
\end{theorem}
% [Tutor Note: 二维情形 MMP 已由本定理完全解决：终点只有三种，且收缩次数受 $\rho(X)-1$ 控制。三维的困难正来自翻转，这也是第 6 章的主题。]

在上述定理中，Q-分解性假设对二维 dlt 对总是成立（命题 4.11），故不必要，这一点应予注意。

\section{极小模型与标准模型}
\takeaway{本节把极小模型推广为相对的 $f$-极小模型，定义弱伪标准模型、伪极小模型与标准模型，并证明它们与 MMP 的终点一致，且在最坏相差余维数 1 的意义下唯一。}

本节的结果在 $X$ 为域上正规概形或正规解析空间、$f:X\to S$ 为固有态射（解析情形为与射影态射双有理同型的固有态射）时成立。对的极小模型定义，是为适配通过 MMP 得到的结论 3.31 阶段 3(2) 而成。定义 2.13 的一个侧面容易一般化。

\begin{definition}{3.48}
设 $(X,\Delta)$ 为对，$f:X\to S$ 为固有态射。当 $(X,\Delta)$ 为 dlt 且 $K_X+\Delta$ 为 $f$-nef 时，$(X,\Delta)$ 称为 $f$-极小（$f$-minimal）或相对极小（relatively minimal）。
\end{definition}
% [原文 p.134]
定义 2.13 还有一个部分，它决定何时 $X$ 为 $X'$ 的极小模型。如 3.30 所见，对的双有理等价概念存在问题，故必须给出「某对成为其他对的极小模型」这一定义。下例表明此时须加注意。

\begin{example}{3.49}
设 $S$ 是只有一个奇点 $P\in S$ 的射影曲面，$f:S'\to S$ 为其极小奇点消解。设 $K_S$ 丰富，$\operatorname{Ex}(f)$ 是不可约非奇异有理曲线 $D$，且 $(D^2)=-n$。此时 $(S,0)$ 为 klt，且
\[
K_{S'}\equiv f^*K_S-\frac{n-2}{n}D.
\]
对 $(S',cD)$ 运行 MMP。若 $c>\frac{n-2}{n}$，最初阶段即收缩 $D$ 并结束。若 $c\le\frac{n-2}{n}$，则 $(D\cdot(K_{S'}+cD))\ge0$，故不收缩 $D$，$S'$ 常常自身就是其极小模型。

如此，$(S,0)$ 是极小的，但未必是 $(S',cD)$ 的极小模型。
\end{example}
% [Tutor Note: 此例说明「极小」不能只按底空间判断：$(S,0)$ 与 $(S',cD)$ 的差异系数不同，后者可由 $D$ 的系数 $c$ 调节。这正是必须比较差异系数的原因。]

记住此例，就会注意到必须比较极小模型的差异系数与原对的差异系数。

\begin{definition}{3.50}
设 $X$ 为正规簇，$\Delta$ 为使 $K_X+\Delta$ 为 Q-Cartier 的（不必有效的）Q-除子，$(X,\Delta)$ 为对，$f:X\to S$ 为固有态射。考虑可换图
\begin{center}
\begin{tikzcd}
X \arrow[rr, dashed, "\varphi"] \arrow[dr, "f"'] & & X^{w} \arrow[dl, "f^{w}"] \\
& S &
\end{tikzcd}
\end{center}
图中出现的对 $(X^{w},\Delta^{w})$ 与双有理映射 $\varphi:X\dashrightarrow X^{w}$，当满足以下条件时称为 $S$ 上 $(X,\Delta)$ 的弱伪标准模型（weak pseudo canonical model）。
\begin{enumerate}[label=(\arabic*)]
\item $f^{w}$ 为固有。
\item $\varphi^{-1}$ 没有例外除子。
\item $\Delta^{w}=\varphi_*\Delta$。
\item $K_{X^{w}}+\Delta^{w}$ 为 Q-Cartier 且为 $f^{w}$-nef。
\item 对各 $\varphi$-例外除子 $E\subset X$，$a(E,X,\Delta)\le a(E,X^{w},\Delta^{w})$。
\end{enumerate}
\end{definition}
% [Tutor Note: 这里对 $X^{w}$ 的奇点没有任何假设，故 (5) 是唯一控制奇点的条件；(2) 排除 $\varphi^{-1}$ 自带例外除子的退化情形，使差异系数可比。]

% [原文 p.135]
弱标准模型 $(X^{m},\Delta^{m})=(X^{w},\Delta^{w})$ 称为 $S$ 上 $(X,\Delta)$ 的伪极小模型（pseudo minimal model），当在上述 (1)--(5) 之外还满足
\begin{enumerate}
\item[（$5^{m}$）] 对各 $\varphi$-例外除子 $E\subset X$，$a(E,X,\Delta)<a(E,X^{m},\Delta^{m})$。
\end{enumerate}
弱标准模型 $(X^{c},\Delta^{c})=(X^{w},\Delta^{w})$ 称为 $S$ 上 $(X,\Delta)$ 的伪标准模型（pseudo canonical model），当在上述 (1)--(5) 之外还满足
\begin{enumerate}
\item[（$4^{c}$）] $K_{X^{c}}+\Delta^{c}$ 为 $f^{c}$-丰富。
\end{enumerate}
但这些概念未对 $(X,\Delta)$ 的奇点作任何假设，故多数情形不存在。$(X,\Delta)$ 为 lc 对且 $\Delta$ 有效时，弱伪标准模型（或伪标准模型）称为弱标准模型（weak canonical model）（或标准模型（canonical model））。又，$(X,\Delta)$ 为 dlt 对时，伪极小模型称为极小模型（minimal model）。表面上看不同，但这里定义的弱标准模型与 [Sho92] 的定义等价。

现从 Q-分解 dlt 对 $(X,\Delta)$ 与射影态射 $f:X\to S$ 出发，运行 $S$ 上的 MMP（3.31）。若纲领可行，则得到对 $(X^{*},\Delta^{*})$ 与射影态射 $f^{*}:X^{*}\to S$。若 MMP 以极小模型情形（3.31 阶段 3(2)）终止，则 $(X^{*},\Delta^{*})$ 是 $(X,\Delta)$ 的极小模型。事实上 (1)--(4) 显然成立，且 $\varphi:X\dashrightarrow X^{*}$ 是除子收缩与翻转的合成，故在每一步用引理 3.38 即得 $(5^{m})$。

在定义 2.13 中，当 $X$ 非奇异时定义了 $X$ 的极小模型 $X'$；其 $(X',0)$ 在上述定义下也是 $(X,0)$ 的极小模型。事实上 (1) 与 (4) 是定义 2.13 的一部分，$\Delta=\Delta'=0$ 使 (3) 显然成立。令 $E$ 为 $\varphi:X\dashrightarrow X'$ 的例外除子，则 $a(E,X,0)=0$，又 $X'$ 为终端（定义 2.13），故 $a(E,X',0)>0$。于是 $(5^{m})$ 成立。另外，若存在 $\varphi^{-1}$-例外除子 $E'$，则因 $X$ 非奇异有 $a(E',X,0)>0$，另一方面由命题 3.51 有 $a(E',X,0)\le a(E',X',0)=0$。矛盾，故 (2) 也成立。\hfill$\square$
% [Tutor Note: 这一段说明定义 2.13 是本节的特例：$\Delta=0$、$X$ 非奇异。注意 (2) 的验证借用了命题 3.51，而后者要到下一页才证明，阅读时宜先跳过回头再看。]

% [原文 p.136]
这些结果表明，它们至少不会使 $(X,\Delta)$ 的奇点程度变坏。

\begin{propositionp}{3.51}
在定义 3.50 中，假设 $(X^{w},\Delta^{w})$ 满足 (1)、(4)、(5) 与下列条件
\begin{enumerate}
\item[（$3'$）] $\varphi_*\Delta-\Delta^{w}$ 去掉 $\varphi^{-1}$-例外除子的部分后有效，
\end{enumerate}
则下列条件满足。
\begin{enumerate}
\item[（$5'$）] $X$ 上方的每个除子 $E$ 满足 $a(E,X,\Delta)\le a(E,X^{w},\Delta^{w})$。
\end{enumerate}
（即用更弱的 $(3')$ 代替 (3)。特别地，$(X,\Delta)$ 的弱伪标准模型 $(X^{w},\Delta^{w})$ 满足 $(5')$。）
\end{propositionp}

\begin{proof}
任取 $X$ 上方的除子 $E$，考虑使 $\operatorname{center}_Y E$ 成为除子的如下交换图。
\begin{center}
\begin{tikzcd}
& Y \arrow[dl, "g"'] \arrow[dr, "h"] & \\
X \arrow[rr, dashed, "\varphi"] \arrow[dr, "f"'] & & X^{w} \arrow[dl, "f^{w}"] \\
& S &
\end{tikzcd}
\end{center}
下列式子是标准除子的比较公式。
\[
K_Y \equiv g^*(K_X+\Delta)+E_1,
\]
\[
K_Y \equiv h^*(K_{X^{w}}+\Delta^{w})+E_2.
\]
注意 $a(E,X^{w},\Delta^{w})-a(E,X,\Delta)$ 是 $E_2-E_1$ 中 $E$ 的系数。记
\[
B:=g^*(K_X+\Delta)-h^*(K_{X^{w}}+\Delta^{w})\equiv E_2-E_1.
\]
于是 $-B$ 为 $g$-nef（定义 3.50(4)），且 $g_*B=g_*(E_2-E_1)$ 由 $(3')$ 与定义 3.50(5) 有效。故由引理 3.39，$E_2-E_1$ 有效，证明完成。
\end{proof}

极小模型与标准模型的存在本质上与 MMP 等价。下面考虑这些模型在何种程度上唯一。

\begin{theoremp}{3.52}
$(X,\Delta)$ 是 $K_X+\Delta$ 为 Q-Cartier 的对，$f:X\to S$ 为固有态射。
\begin{enumerate}[label=(\arabic*)]
\item 伪标准模型 $(X^{c},\Delta^{c})$ 唯一，且
% [原文 p.137]
\[
X^{c}=\operatorname{Proj}_S \bigoplus_{m\ge0} f_*\mathcal{O}_X(mK_X+\lfloor m\Delta\rfloor).
\]
\item $(X,\Delta)$ 的任意两个伪极小模型在余维数 1 处同型。
\end{enumerate}
\end{theoremp}

\begin{proof}
使用命题 3.51 的证明中所用的图。选取 $\Delta_Y$，使得 $K_Y+\Delta_Y=g^*(K_X+\Delta)$ 且 $g_*\Delta_Y=\Delta$；则由命题 3.51，对弱伪标准模型 $(X^{w},\Delta^{w})$ 的任意射 $h:Y\to X^{w}$，可得下式：
\[
K_Y+\Delta_Y \equiv h^*(K_{X^{w}}+\Delta^{w})+(\text{有效 } h\text{-例外除子}).
\]
先处理 $(X^{w},\Delta^{w})=(X^{c},\Delta^{c})$ 的情形。此时
\[
(X^{w},\Delta^{w})=\operatorname{Proj}_S \bigoplus_{m\ge0} f^{w}_*\mathcal{O}_{X^{w}}(mK_{X^{w}}+\lfloor m\Delta^{w}\rfloor)
\]
\[
=\operatorname{Proj}_S \bigoplus_{m\ge0} (f\circ g)_*\mathcal{O}_Y(mK_Y+\lfloor m\Delta_Y\rfloor)
\]
\[
=\operatorname{Proj}_S \bigoplus_{m\ge0} f_*\mathcal{O}_X(mK_X+\lfloor m\Delta\rfloor).
\]
这意味着伪标准模型唯一。

其次，设 $\varphi_i:X\dashrightarrow X_i^{m}$ 为两个伪极小模型。须证明 $\varphi_1$ 与 $\varphi_2$ 具有相同的例外除子。选取固有双有理态射 $g:Y\to X$ 以及由 $g$ 诱导的双有理映射 $h_i:Y\to X_i^{m}$，使两者都是射。在标准除子公式
\[
K_Y+\Delta_Y \equiv h_i^*(K_{X_i^{m}}+\Delta_i^{m})+Z_i
\]
中，$Z_i$ 是有效 $h_i$-例外除子，这在证明开头已确认。又由定义 3.50(3) 与定义 3.50$(5^{m})$ 可知，$\operatorname{Supp}g_*Z_i$ 的除子部分与 $\varphi_i$ 的全部例外除子之并一致。由两个公式得到
\[
h_1^*(K_{X_1^{m}}+\Delta_1^{m})-h_2^*(K_{X_2^{m}}+\Delta_2^{m})\equiv Z_2-Z_1.
\]
对 $h_1:Y\to X_1^{m}$（或 $h_2:Y\to X_2^{m}$）应用引理 3.39，可得 $Z_2-Z_1$（或 $Z_1-Z_2$）有效。于是 $Z_1=Z_2$，$\varphi_1$ 与 $\varphi_2$ 具有相同的例外除子。
\end{proof}

\begin{corollaryp}{3.53}
设 $(X,\Delta)$ 与 $(X',\Delta')$ 是使 $K_X+\Delta$ 与 $K_{X'}+\Delta'$ 为 Q-Cartier 的对，$p:X'\to X$ 与 $f:X\to S$ 为固有态射。假设 $K_{X'}+\Delta'-p^*(K_X+\Delta)$ 有效且为 $p$-例外。
% [原文 p.138]
则 $(X,\Delta)$ 与 $(X',\Delta')$ 具有 $S$ 上同型的伪标准模型。
\end{corollaryp}

\begin{proof}
由假设，$p_*\mathcal{O}_{X'}(mK_{X'}+\lfloor m\Delta'\rfloor)=\mathcal{O}_X(mK_X+\lfloor m\Delta\rfloor)$ 成立。因此
\[
(f\circ p)_*\mathcal{O}_{X'}(mK_{X'}+\lfloor m\Delta'\rfloor)=f_*\mathcal{O}_X(mK_X+\lfloor m\Delta\rfloor).
\]
应用定理 3.52 即可。
\end{proof}

\begin{corollaryp}{3.54}
设 $f_i:X_i\to S$（$i=1,2$）为固有态射，$X_i$ 为终端，且 $K_{X_i}$ 为 $f_i$-nef。此时 $S$ 上任意双有理映射 $\varphi:X_1\dashrightarrow X_2$ 在余维数 1 处为同构。
\end{corollaryp}

\begin{proof}
设 $g_i:X'\to X_i$ 为公共奇点消解。此时对 $i=1,2$，$X_i$ 是 $X'$ 的极小模型，这在定义 3.50 末尾已确认。故由定理 3.52，$\varphi$ 在余维数 1 处为同构。
\end{proof}

\chapterreview{本章是全书的技术核心，其任务是证明锥定理，并把曲线锥理论推进到可直接驱动极小模型纲领（MMP）的程度。第 1 章在曲面与非奇异情形给出雏形，第 2 章准备了奇点、消没定理与相对化的语言；本章则把这些工具组装成一般维数的锥定理、无基点定理与有理性定理，使 MMP 的每一步收缩与翻转都能在除子层面得到控制。第 3 章回答的是「MMP 为什么能运行」，而剩下的「翻转为什么存在」则交给第 6 章。

从脉络看，§3.1 引入 Fano 簇与 dlt 对，§3.2 与 §3.3 分别给出无基点定理与锥定理，§3.4 与 §3.5 是对应的有理性定理与非消没定理，§3.6 把它们相对化，§3.7 由此推出 Q-分解性与奇点类别在收缩、翻转下不变，§3.8 再把极小模型推广为相对的 $f$-极小模型与（弱）伪标准模型。后三章反复使用的正是锥定理与相对化版本，故本章每个条件都值得细读。

核心结论可列成一条链：锥定理（定理 3.25，其 dlt 版为定理 3.35）给出曲线锥的可数分解；无基点定理（定理 3.3）把丰富性化为数值条件；有理性定理保证 $\R$-除子的有理结构；非消没定理给出区间的正性；相对锥定理把上述结论搬到 $f:X\to S$ 上。由此得到命题 3.36、命题 3.37、推论 3.42--3.44：除子收缩使 $\rho$ 降 1，翻转保持 $\rho$ 与奇点类别。3.45 与 3.46 则诚实指出终止性与翻转存在性仍是未解问题。

辨析与提醒：其一，本书尾符有三档，「证明紧跟」的声明末尾无符号，不可与「书中不证」（■）或「证明在后文」（□）混淆。其二，「极射线」与「极端的／射线型」是两个概念，后者只是相对 Picard 秩 1 的条件。其三，翻转收缩是小收缩，且定义中的 $D$ 不必有效，故定义可直接沿 MMP 使用。其四，锥定理的 dlt 版证明采用扰动技巧，把 $\Delta$ 换成 $(1-\delta)\Delta+\eps H$ 以化为 klt，这是最值得记住的技巧。

评注：本章内容取自 Kollár--Mori 的原始论文与 1990 年代的综合工作，定理 3.35 的相对锥定理与三维翻转的存在性（森重文，1988）共同支撑了三维 MMP。定理 3.47 表明二维 MMP 已完全解决，而高维的情形直到 2000 年代的极小模型纲领工作才取得突破。阅读时宜把 §3.3 与 §3.6 对照：先理解绝对锥定理的几何含义，再看相对化如何把「存在丰富除子」换成「存在 $f$-丰富除子」；引理 3.38 与 3.39 是其后所有唯一性论证的公共工具。}
